% !TeX program = xelatex
% =============================================================
%  甘肃农业大学实验报告模板
%  基于 XeLaTeX
%
%  文档结构：
%    第 1 页 ............ 封面
%    第 2 页 ............ 信息栏 + 实验正文（带“装订线/批阅意见”边框）
%    第 3 页及以后 ...... 与第 2 页同版式，用于续写
%
%  文件说明：
%    preamble.tex                            宏包与版面样式
%    cover.tex                               封面
%    information.tex                         关键信息
%    sections/*.tex                          实验正文章节
%    media /The_format_requires              图像资源
% =============================================================
\documentclass[UTF8,a4paper,zihao=-4,fontset=fandol]{ctexart}
% 修改对应部分就可以进行整个文档编辑

% 注意：下面每个 \newcommand 定义的都是“字符串内容”，改等号后面花括号里的东西即可，不需要则留空

% 成绩
\newcommand{\infoScore}{}




%=========================================================================
% 
%                       课    程    信    息
%
%=========================================================================
% 实验题目
\newcommand{\infoExpNo}{实验N}
\newcommand{\infoExpTitle}{实验名}
% 课程名称
\newcommand{\infoCourse}{课程名}
% 实验日期（格式 2026-09-01）
\newcommand{\infoExperimentTime}{实验日期}
% 任课教师姓名
\newcommand{\infoTeacher}{教师名}
% 总学时
\newcommand{\infoHours}{学时}
% 小组编号
\newcommand{\infoGroup}{编号}




%=========================================================================
% 
%                       个    人    信    息
%
%=========================================================================

% 学生姓名
\newcommand{\infoName}{李四}
% 学号
\newcommand{\infoStudentId}{20240001}
% 学院名称
\newcommand{\infoCollege}{信息科学技术学院}
% 班级名称
\newcommand{\infoClass}{aaa}
% 班级人数
\newcommand{\infoStudentCount}{60}
% 学年开始年份
\newcommand{\infoYearStart}{26}
% 学年结束年份
\newcommand{\infoYearEnd}{27}
% 学期
\newcommand{\infoSemester}{1}
% 专业名称
\newcommand{\infoMajor}{aaa}
% 年级（24 = 2024 级）
\newcommand{\infoGrade}{24}
% 班级序号（2 = 2 班）
\newcommand{\infoClassNo}{2}


% =============================================================
%  preamble.tex  ——  宏包、版面参数与页面样式
%  编译引擎：XeLaTeX（OverLeaf 请在编译器中选 XeLaTeX）
% =============================================================

\usepackage{geometry}
\usepackage{graphicx}
\usepackage{eso-pic}
\usepackage{tikz}
\usetikzlibrary{calc}
\usepackage{enumitem}
\usepackage{array}

% 带圈数字 ①②③… 归入中文字体（否则落到拉丁字体 Latin Modern，会缺字）
\xeCJKDeclareCharClass{CJK}{"2460 -> "24FF}

% ---------- 版面参数 ----------
% 封面
\newcommand{\CoverPaperGeometry}{%
  \newgeometry{left=2.5cm,right=2.5cm,top=2.5cm,bottom=2.5cm}}
% 报告页（第 2 页起）：左边界＝边框横线左端（模版实测 2.61cm）；
% 右边界 4.3cm → 正文右边界≈16.7cm，位于批阅意见竖线（17.27cm）左侧留 0.5cm 间距
\newcommand{\ReportPaperGeometry}{%
  \newgeometry{left=2.61cm,right=4.3cm,top=4.40cm,bottom=2.2cm}}

% ---------- 报告页边框坐标（单位 cm，相对页面绝对定位；数值取自模版实测）----------
\def\FRleft{1.31}    % “装订线”虚线横坐标
\def\FRbody{2.61}    % 横线左端（与信息栏/正文左边界平齐）
\def\FRright{17.27}  % “批阅意见”栏竖线横坐标
\def\FRlineR{19.55}  % 横线右端（向右盖过批阅意见栏）
\def\FRbot{1.18}     % 下双线·粗线 距页底（＝距页顶 28.52cm）
\def\FRbott{1.08}    % 下双线·细线 距页底（＝距页顶 28.62cm）
\def\LWthick{2.6pt}  % 粗线线宽（≈0.9mm）
\def\LWthin{0.8pt}   % 细线线宽（≈0.28mm）

% 绝对页面坐标（相对页面左上/左下角）。用 calc 库写法，
% 宏参数可正常展开（[xshift=...] 前缀在宏展开时 tikz 无法识别）。
\newcommand{\PtNW}[2]{($(current page.north west)+(#1cm,-#2cm)$)}
\newcommand{\PtSW}[2]{($(current page.south west)+(#1cm,#2cm)$)}

% ---------- 报告页 页眉 + 边框（背景层，报告页逐页绘制）----------
% 报告页页数计数器：在 shipout 时递增，用于区分“第2页（带信息栏）”与后续续写页
\newcount\ReportPageNo

\newcommand{\ReportPageArt}{%
  \global\advance\ReportPageNo by 1
  \begingroup
  % 第 2 页顶部横线位于信息栏下方；后续续写页位于页眉正下方
  \ifnum\ReportPageNo=1 \def\RuleTop{7.08}\else \def\RuleTop{4.20}\fi
  % 上双线：粗线中心 / 细线中心
  \pgfmathsetmacro{\RuleTap}{\RuleTop+0.05}
  \pgfmathsetmacro{\RuleTbp}{\RuleTop+0.15}
  % “批阅意见：”标签高度
  \pgfmathsetmacro{\LabelY}{\RuleTop+0.55}
  % “装订线”三字在框内垂直居中、字距 1.9cm；上下两段虚线的端点
  \pgfmathsetmacro{\BindC}{(\RuleTop+28.52)/2}
  \pgfmathsetmacro{\BindA}{\BindC-2.35}
  \pgfmathsetmacro{\BindB}{\BindC+2.35}
  \pgfmathsetmacro{\BindZz}{\BindC-1.9}
  \pgfmathsetmacro{\BindZl}{\BindC+1.9}
  \begin{tikzpicture}[remember picture,overlay]
    % 页眉：校徽校名
    \node[anchor=north west,inner sep=0pt] at \PtNW{2.70}{1.70}
      {\includegraphics[width=7.0cm]{media/The_format_requires/gyau-logo.png}};
   % 页眉：________学院实验（实习）报告
\node[anchor=north east,inner sep=0pt,font=\zihao{4}] at \PtNW{19.15}{3.35}
  {\underline{\makebox[4.3cm][c]{\infoCollege}}学院实验（实习）报告};
    % 上双线（上粗下细）
    \draw[line width=\LWthick] \PtNW{\FRbody}{\RuleTap} -- \PtNW{\FRlineR}{\RuleTap};
    \draw[line width=\LWthin]  \PtNW{\FRbody}{\RuleTbp} -- \PtNW{\FRlineR}{\RuleTbp};
    % 下双线（上粗下细）
    \draw[line width=\LWthick] \PtSW{\FRbody}{\FRbot} -- \PtSW{\FRlineR}{\FRbot};
    \draw[line width=\LWthin]  \PtSW{\FRbody}{\FRbott} -- \PtSW{\FRlineR}{\FRbott};
    % 装订线·上段虚线（顶双线 → “装”字之前）
    \draw[dashed,line width=0.6pt] \PtNW{\FRleft}{\RuleTop} -- \PtNW{\FRleft}{\BindA};
    % 装订线·下段虚线（“线”字之后 → 底双线），末端向右拐与底部粗线相接
    \draw[dashed,line width=0.6pt] \PtNW{\FRleft}{\BindB} -- \PtSW{\FRleft}{\FRbot}
      -- \PtSW{\FRbody}{\FRbot};
    % 装 订 线（三字，压在虚线上）
    \node[anchor=center,inner sep=0pt] at \PtNW{\FRleft}{\BindZz} {装};
    \node[anchor=center,inner sep=0pt] at \PtNW{\FRleft}{\BindC}  {订};
    \node[anchor=center,inner sep=0pt] at \PtNW{\FRleft}{\BindZl} {线};
    % 批阅意见栏竖线
    \draw[line width=0.7pt] \PtNW{\FRright}{\RuleTap} -- \PtSW{\FRright}{\FRbot};
    % 批阅意见标签
    \node[anchor=north west,inner sep=0pt,font=\footnotesize]
      at \PtNW{17.35}{\LabelY} {批阅意见：};
  \end{tikzpicture}%
  \endgroup}

\newif\ifReportPage
\AddToShipoutPictureBG{\ifReportPage\ReportPageArt\fi}

% ---------- 额外设定 ----------


% 实验报告小节标题（如“一、实验目的”）
\newcommand{\ExpSection}[1]{%
  \par\addvspace{0.85em}%
  \noindent{\zihao{4}\bfseries #1}\par\nobreak\addvspace{0.35em}}

\linespread{1.5}   % 正文行距

\pagestyle{empty}  % 不显示页码


\begin{document}

% ==================== 第 1 页：封面 ====================
% =============================================================
%  cover.tex  ——  封面部分
% =============================================================
\CoverPaperGeometry

\vspace*{0.40cm}
\hspace*{0.45cm}\includegraphics[width=7.2cm]{media/The_format_requires/gyau-logo.png}\par

\vspace{3.4cm}
\begin{center}
  {\zihao{-0}\bfseries 本科生实验报告}
\end{center}

\vspace{2.2cm}
\begin{center}
\zihao{4}\linespread{1.0}\selectfont
\begin{tabular}{@{}l@{}}
学\hspace{1.3em}院\hspace{0.15em}：\underline{\makebox[6.0cm][c]{\infoCollege}}\\[1.15em]
课程名称：\underline{\makebox[6.0cm][c]{\infoCourse}}\\[1.15em]
班\hspace{1.3em}级\hspace{0.15em}：\underline{\makebox[3.4cm][c]{\infoClass}}\quad 人数：\underline{\makebox[1.2cm][c]{\infoStudentCount}}\\[1.15em]
学年学期：20\underline{\makebox[0.8cm][c]{\infoYearStart}}——20\underline{\makebox[0.8cm][c]{\infoYearEnd}}学年第\underline{\makebox[0.8cm][c]{\infoSemester}}学期\\[1.15em]
学\hspace{0.6em}时\hspace{0.6em}数\hspace{0.15em}：\underline{\makebox[6.0cm][c]{\infoHours}}\\[1.15em]
指导教师：\underline{\makebox[6.0cm][c]{\infoTeacher}}\\
\end{tabular}
\end{center}

\vfill
\begin{center}
  \zihao{4}教\hspace{0.6em}务\hspace{0.6em}处\hspace{0.6em}制
\end{center}
\vspace*{0.6cm}


% ==================== 第 2 页起：报告正文 ====================
\newpage
\ReportPaperGeometry
\ReportPagetrue          % 打开报告页边框




% ---- 信息栏----
% 三行铺满 x=2.61~19.37cm
\begingroup
\zihao{5}\linespread{1.3}\selectfont
\def\RowW{16.76cm}
\noindent\makebox[0pt][l]{\hbox to \RowW{%
  \makebox[9.0cm][l]{专业班级：\underline{\makebox[1.8cm][c]{\infoMajor}}专业\underline{\makebox[0.9cm][c]{\infoGrade}}级\underline{\makebox[0.9cm][c]{\infoClassNo}}班}%
  组别：\underline{\makebox[1.4cm][c]{\infoGroup}}\hfill 指导教师：\underline{\makebox[2.2cm][c]{\infoTeacher}}\hss}}\par
\vspace{0.31cm}
\noindent\makebox[0pt][l]{\hbox to \RowW{%
  \makebox[9.0cm][l]{课程名称：\underline{\makebox[5.0cm][c]{\infoCourse}}}%
  姓名：\underline{\makebox[1.4cm][c]{\infoName}}\hfill 学\hspace{0.4em}号：\underline{\makebox[2.0cm][c]{\infoStudentId}}\hss}}\par
\vspace{0.31cm}
\noindent\makebox[0pt][l]{\hbox to \RowW{%
  \makebox[9.0cm][l]{实验（实习）时间：\underline{\makebox[4.2cm][c]{\infoExperimentTime}}}%
  \hfill 成\hspace{0.4em}绩：\underline{\makebox[2.0cm][c]{\infoScore}}\hss}}\par
\endgroup

\vspace*{0.55cm}          % 使正文从第 2 页顶部双线下方开始

% ---- 实验题目 ----
\begin{center}
  {\zihao{4}\bfseries \infoExpNo\quad \infoExpTitle}
\end{center}

% ---- 正文各章节 ----
% ===== 一、实验目的 =====
\ExpSection{一、实验目的}

%补充实际内容即可
% ===== 二、实验仪器设备 =====
\ExpSection{二、实验仪器设备}

%补充实际内容即可

% ===== 三、实验内容及步骤 =====
\ExpSection{三、实验内容及步骤}

%补充实际内容即可
% ===== 四、实验过程 =====
\ExpSection{四、实验过程}

%补充实际内容即可



% ===== 五、实验总结 =====
\ExpSection{五、实验总结}

%补充实际内容即可

\end{document}
